\section{Method}
\label{sec:method}

We present AnomaClaw, a training-free system that combines dense patch-level expert analysis with VLM reasoning for cross-domain visual anomaly detection.
The system operates in three stages---visual retrieval, patch expert analysis, and expert-informed VLM reasoning---with an optional family-adaptive fusion step that calibrates the expert--VLM balance per domain type.

\subsection{Problem Formulation}
\label{sec:formulation}

Given a query image $\query$, a domain code $\domcode$ identifying the inspection context (e.g., \texttt{mvtec\_metal\_nut}, \texttt{brain\_mri}), and a bank of normal images $\bank_\domcode = \{I^n_1, \ldots, I^n_N\}$, the goal is to produce an anomaly score $\ascore \in [0,1]$ where higher values indicate greater anomaly likelihood.
The system operates in a \emph{training-free, few-shot} setting: no parameters are learned from anomaly examples, and only $k$ normal references (typically $k \leq 5$) are used per query at inference time.

We organize the benchmark domains into four \emph{anomaly families} that reflect the nature of deviations each domain exhibits. Each domain code $\domcode$ is assigned a known family label $\family(\domcode)$:
\begin{equation}
    \family(\domcode) \in \{\texttt{local\_appearance},\; \texttt{semantic\_medical},\; \texttt{semantic\_scene},\; \texttt{logical\_structural}\}.
    \label{eq:families}
\end{equation}
Local appearance anomalies (industrial surfaces, retail products) manifest as spatially localized texture or color deviations.
Semantic medical anomalies (brain MRI, liver CT, endoscopy) require clinical interpretation that goes beyond pixel-level comparison.
Semantic scene anomalies (road hazards) involve contextual reasoning about object placement and safety.
Logical structural anomalies (assembly errors, functional defects) concern violations of compositional rules where individual components appear normal in isolation.
This family structure is used for the optional fusion step (\Cref{sec:fusion}) and for selecting family-appropriate guidance text in the VLM prompt; the pipeline \emph{architecture} (retrieval $\to$ patch expert $\to$ VLM) is identical across all domains, but the textual domain knowledge and family-specific guidance provided to the VLM differ per domain.

\subsection{Pipeline Overview}
\label{sec:overview}

\Cref{fig:overview} illustrates the three-stage pipeline.
\textbf{Stage~1 (Visual Retrieval):} We embed all normal images in $\bank_\domcode$ using DINOv2 and retrieve the $k$ most similar references to the query via cosine similarity.
\textbf{Stage~2 (Patch Expert Analysis):} A PatchCore-style module extracts dense patch tokens from the query and references, computes patch-level nearest-neighbor distances, and aggregates these into both a scalar anomaly score $\ascore_{\text{expert}}$ and a structured text summary.
\textbf{Stage~3 (VLM Reasoning):} The VLM receives the query image, retrieved reference images, domain-specific knowledge, and the expert's text summary, then outputs a final anomaly judgment.
An optional \textbf{Stage~4 (Family-Adaptive Fusion)} combines $\ascore_{\text{expert}}$ and the VLM score $\ascore_{\text{vlm}}$ using family-specific weights.

\begin{algorithm}[t]
\caption{AnomaClaw inference pipeline}
\label{alg:pipeline}
\begin{algorithmic}[1]
\REQUIRE Query image $\query$, domain code $\domcode$, normal bank $\bank_\domcode$, number of references $k$
\ENSURE Anomaly score $\ascore \in [0,1]$
\STATE \textbf{// Stage 1: Visual Retrieval}
\STATE $\feat^q \gets \textsc{DINOv2-CLS}(\query)$
\STATE $\refset \gets \text{top-}k \text{ images from } \bank_\domcode \text{ by } \cos(\feat^q, \feat^n)$
\STATE \textbf{// Stage 2: Patch Expert Analysis}
\STATE $\{\patch^q_i\}_{i=1}^{M} \gets \textsc{DINOv2-Patch}(\query)$
\STATE $\mathcal{P}_{\text{ref}} \gets \bigcup_{I^r \in \refset} \textsc{DINOv2-Patch}(I^r)$
\STATE $d_i \gets 1 - \max_{\patch \in \mathcal{P}_{\text{ref}}} \cos(\patch^q_i, \patch)$ for each $i$
\STATE $\ascore_{\text{expert}} \gets \text{mean of top-}1\%\text{ of } \{d_i\}_{i=1}^{M}$
\STATE $t_{\text{expert}} \gets \textsc{FormatSummary}(\ascore_{\text{expert}}, \cos(\feat^q, \feat^{r_1}))$
\STATE \textbf{// Stage 3: VLM Reasoning}
\STATE $\ascore_{\text{vlm}} \gets \textsc{VLM}(\query, \refset, t_{\text{expert}}, t_{\text{domain}})$
\STATE \textbf{// Stage 4: Family-Adaptive Fusion (optional)}
    \STATE $\ascore \gets \alphafam_{\family(\domcode)} \cdot \ascore_{\text{expert}} + (1 - \alphafam_{\family(\domcode)}) \cdot \ascore_{\text{vlm}}$
\RETURN $\ascore$
\end{algorithmic}
\end{algorithm}

\subsection{Visual Retrieval}
\label{sec:retrieval}

Effective reference selection is critical for few-shot anomaly detection: the VLM and patch expert both benefit from references that are visually close to the query, as these reduce the false-positive rate from irrelevant comparisons.
We use DINOv2-Small (ViT-S/14)~\citep{oquab2024dinov2} as the retrieval backbone.
Each image in the normal bank $\bank_\domcode$ is passed through DINOv2, and the CLS token embedding $\feat^n \in \R^{384}$ is $\ell_2$-normalized and stored.
At inference, the query embedding $\feat^q$ is computed identically, and the top-$k$ normal images are retrieved by cosine similarity:
\begin{equation}
    \refset = \underset{S \subset \bank_\domcode,\, |S|=k}{\argmax} \sum_{I^n \in S} \cos(\feat^q, \feat^n).
    \label{eq:retrieval}
\end{equation}
This retrieval step replaces the random reference selection used in prior VLM-based methods~\citep{jiang2024mmad,zhang2025agentiad}, which we find to be particularly harmful in medical and industrial domains where within-class appearance variation is high (e.g., different brain slices or different metal nut orientations).
The retrieved reference set $\refset$ is used by both the patch expert and the VLM in subsequent stages.

\subsection{Patch Expert Analysis}
\label{sec:expert}

The patch expert provides a quantitative grounding signal that the VLM cannot compute from pixels alone.
We extract dense patch tokens from DINOv2 at its native resolution of $518 \times 518$ pixels, yielding a $37 \times 37$ grid of $384$-dimensional feature vectors per image ($M = 1{,}369$ patches).

\paragraph{Reference patch bank.}
For the $k$ retrieved reference images, we concatenate all patch tokens into a reference bank:
\begin{equation}
    \mathcal{P}_{\text{ref}} = \bigcup_{I^r \in \refset} \{\patch^r_1, \ldots, \patch^r_M\}, \quad |\mathcal{P}_{\text{ref}}| = k \cdot M.
    \label{eq:patchbank}
\end{equation}

\paragraph{Patch-level anomaly scoring.}
For each query patch token $\patch^q_i$, we compute the cosine distance to its nearest neighbor in $\mathcal{P}_{\text{ref}}$:
\begin{equation}
    d_i = 1 - \max_{\patch \in \mathcal{P}_{\text{ref}}} \cos(\patch^q_i, \patch).
    \label{eq:patch_distance}
\end{equation}
The expert anomaly score is the mean of the top-$1\%$ most anomalous patch distances:
\begin{equation}
    \ascore_{\text{expert}} = \frac{1}{|\mathcal{T}|} \sum_{i \in \mathcal{T}} d_i, \quad \mathcal{T} = \{i : d_i \geq d_{(0.99M)}\},
    \label{eq:expert_score}
\end{equation}
where $d_{(0.99M)}$ denotes the $99$th-percentile distance.
This aggregation captures localized deviations---characteristic of surface defects, lesions, and structural misalignments---while remaining robust to the natural patch-level variation present in normal images.

\paragraph{Structured text interpretation.}
Rather than passing the raw scalar score directly for fusion or routing, we convert the expert's findings into a structured text summary $t_{\text{expert}}$ that the VLM can parse:
\begin{quote}
\small
\texttt{Patch-level anomaly score: 0.248 (higher=more different from refs)} \\
\texttt{Global similarity to best reference: 0.864} \\
\texttt{Assessment: moderate anomaly signal}
\end{quote}
The numeric value in the prompt is the raw mean patch distance $\ascore_{\text{expert}}$ (Eq.~\ref{eq:expert_score}).
To generate the qualitative assessment label, this raw distance is internally passed through a sigmoid calibration $\sigma(20 \cdot (d_{\text{raw}} - 0.18))$ and binned into five levels (``very similar to normal'' through ``strong anomaly signal'').
The downstream AUROC computation uses the sigmoid-calibrated value; the VLM sees the raw distance plus the qualitative interpretation.
This text summary is the sole channel through which expert evidence reaches the VLM.

\subsection{Expert-Informed VLM Reasoning}
\label{sec:vlm}

The VLM receives a multi-modal prompt consisting of four components:
\begin{enumerate}[leftmargin=*,itemsep=1pt]
    \item \textbf{Reference images:} the $k$ retrieved normal images from $\refset$, establishing the visual baseline.
    \item \textbf{Query image:} the image $\query$ to be classified.
    \item \textbf{Domain knowledge:} a text description specifying the anomaly criteria for domain $\domcode$ and listing common false-positive patterns (e.g., ``normal liver CT may show benign cysts; do not flag these as anomalies'').
    \item \textbf{Expert analysis text:} the structured summary $t_{\text{expert}}$ from the patch expert.
\end{enumerate}
The VLM (GPT-5.4) outputs a JSON response containing an image-level label $y \in \{\texttt{normal}, \texttt{anomalous}\}$ and a confidence value $c \in [0,1]$.
The VLM anomaly score is computed as:
\begin{equation}
    \ascore_{\text{vlm}} =
    \begin{cases}
        c & \text{if } y = \texttt{anomalous}, \\
        1 - c & \text{if } y = \texttt{normal}.
    \end{cases}
    \label{eq:vlm_score}
\end{equation}

\paragraph{Why expert-as-context outperforms alternatives.}
The central design choice in AnomaClaw is to provide expert evidence as \emph{text context} rather than using the expert as a \emph{router} or performing \emph{score fusion}.
We briefly motivate this here; full ablations appear in \Cref{sec:experiments}.

\textbf{Expert-as-router} uses the expert score to decide whether to invoke the VLM at all: if the expert is confident, its score is returned directly, and the VLM is called only for ambiguous cases.
This underperforms in our cross-domain setting because medical images often yield moderately elevated expert scores even when they are normal.
A router can therefore suppress VLM consultation on exactly the cases where semantic reasoning is most needed.

\textbf{Score fusion} averages the expert and VLM scores: $\ascore = \frac{1}{2}(\ascore_{\text{expert}} + \ascore_{\text{vlm}})$.
The fundamental issue is that arithmetic averaging cannot express the conditional statement ``this expert evidence is unreliable for this sample.''
When the expert and VLM disagree sharply, their average can remain mediocre even if one of the two subsystems is clearly correct.

\textbf{Expert-as-context} avoids both failure modes.
When the VLM reads ``Patch score: 0.248, Global similarity: 0.864, Assessment: moderate anomaly signal,'' it can reason: the patch expert detects some deviation, but the high global similarity suggests the query is close to the normal distribution.
The VLM then examines the images directly and makes an informed judgment, discounting the expert when the visual evidence contradicts the quantitative signal.
This design preserves the VLM's full decision authority while grounding it with evidence it could not compute independently.

\subsection{Family-Adaptive Fusion}
\label{sec:fusion}

While expert-as-context is the primary integration mechanism, we observe that the \emph{scalar} expert score $\ascore_{\text{expert}}$ carries complementary information whose reliability varies systematically across anomaly families.
Family-adaptive fusion combines the expert and VLM scores with family-specific weights:
\begin{equation}
    \ascore = \alphafam_{\family(\domcode)} \cdot \ascore_{\text{expert}} + (1 - \alphafam_{\family(\domcode)}) \cdot \ascore_{\text{vlm}},
    \label{eq:fusion}
\end{equation}
where $\alphafam_{\family(\domcode)} \in [0,1]$ is a scalar weight for the anomaly family assigned to domain $\domcode$.

\paragraph{Calibration procedure.}
The fusion weights are tuned on a disjoint calibration pool of 200 items: 20 per domain (10 normal, 10 anomalous), separate from the 1{,}200-item test benchmark.
For each anomaly family, we pool the calibration items from the domains in that family and perform a grid search over $\alphafam_\family \in \{0.0, 0.05, 0.10, \ldots, 1.0\}$, selecting the value that maximizes AUROC on that family's pooled calibration subset.
The resulting weights are frozen for all test evaluation.

\paragraph{Calibrated weights and interpretation.}
The tuned values reveal a clear pattern in expert reliability:
\begin{itemize}[leftmargin=*,itemsep=1pt]
    \item \textbf{Local appearance} ($\alphafam = 0.35$): The patch expert is effective at detecting surface defects and texture anomalies, contributing meaningfully alongside the VLM.
    \item \textbf{Semantic medical} ($\alphafam = 0.05$): The expert is unreliable on medical imagery---subtle pathology often falls within the normal feature distribution---so the VLM dominates.
    \item \textbf{Semantic scene} ($\alphafam = 0.00$): Scene-level anomalies are entirely semantic (e.g., a car driving on the wrong side), and the expert provides no useful signal.
    \item \textbf{Logical structural} ($\alphafam = 0.60$): While the expert cannot reason about logical correctness per se, misplaced or missing components produce detectable appearance differences in the patch feature space. The high calibration-set weight reflects this; however, the test-set AUROC for D9 is only marginally affected by fusion (0.792 vs.\ 0.796 without), so this weight should be interpreted cautiously.
\end{itemize}
With these calibrated weights, AnomaClaw achieves its best operating point of $0.882$ macro AUROC across all 10 domains.
The fusion adds no additional model parameters and negligible computation---it is a single weighted average applied after both scores are already computed.
